% Fuentes: controles EH y cosmológico de la sección de Lie.
\begin{frame}[label=nav-frame-022]{Casos de referencia: momentos y contracción}
\TablaSetup
\begin{tabular}{@{}p{.24\textwidth}p{.44\textwidth}p{.25\textwidth}@{}}
\TablaCabecera{Objeto} & \TablaCabecera{Einstein--Hilbert} & \TablaCabecera{Cosmológico} \\
\TablaLinea
Lagrangiano & $L_{\mathrm{EH}}=R$ & $L_\Lambda=-2\Lambda$ \\[2mm]
Momento de curvatura & $P_{\mathrm{EH}}^{abcd}=\tfrac12(g^{ac}g^{bd}-g^{ad}g^{bc})$ & $P_\Lambda^{abcd}=0$ \\[3mm]
Contracción con Riemann & $\mathcal R_{ab}^{(\mathrm{EH})}=R_{ab}$ & $\mathcal R_{ab}^{(\Lambda)}=0$ \\[2mm]
Momento métrico\newline respecto de $g_{ab}$ & $P^{ab}=-2\mathcal R^{ab}$\newline aplicado al caso EH & $P_\Lambda^{ab}=0$ \\
\TablaLinea
\end{tabular}\par
\vspace{4mm}

La identidad de Lie relaciona el momento métrico con la contracción
de curvatura. El término cosmológico anula ambos momentos.
\end{frame}
